In random-ReLU-features regression with $n=200$ samples, minimum-norm interpolation peaks at $N/n=1$ in every run, and the peak is higher at higher label noise, but it is large already without noise. A fixed ridge lowers the peak by orders of magnitude down to a minimum near $\lambda=10^{-3}$--$10^{-2}$ and rises again for larger $\lambda$, so monotone decrease in the ridge is false. A ridge selected by validation removes the peak to within a few percent of the best error on average on all eight tasks, but not exactly on every seed, and with 5 seeds our registered paired test against interpolation on that hump statistic is not significant even though the hump is orders of magnitude smaller. A single global ridge did as well as per-width tuning here. These results hold for this small finite-size setting and should be checked at larger scale.
